\chapter{Algebraic quantum field theory}\label{ch:aqft}

The Haag--Kastler axioms, the Reeh--Schlieder and Bisognano--Wichmann theorems and the Unruh effect, why local algebras are of type III$_1$, and superselection sectors.

\section{Haag--Kastler axioms}\label{ch:haagkastler}

\begin{physmotiv}
Quantum field theory is, at its core, the statement that observables can be \emph{localized}: associated with regions of spacetime, with those of spacelike separated regions commuting, and with Poincar\'e transformations moving regions around. The Haag--Kastler framework keeps exactly this and nothing more. Fields, Lagrangians, and even particle interpretations are derived or optional. It is the continuum, relativistic version of the quasi-local algebra of \cref{ch:quasilocal}, and it is the setting in which all statements about type III algebras of regions are made.
\end{physmotiv}

\subsection{The axioms}

A \emph{local net} assigns to each open bounded region $O\subset\R^{d}$ (typically double cones) a von Neumann algebra $\A(O)$ on a Hilbert space $\HH$ such that:
\begin{enumerate}[label=(HK\arabic*)]
\item \textbf{Isotony.} $O_1\subset O_2\Rightarrow\A(O_1)\subset\A(O_2)$.
\item \textbf{Locality (Einstein causality).} If $O_1$ and $O_2$ are spacelike separated, $[\A(O_1),\A(O_2)]=0$; equivalently $\A(O')\subset\A(O)'$ with $O'$ the causal complement. (For fermions, graded commutation.)
\item \textbf{Covariance.} There is a strongly continuous unitary representation $U$ of the Poincar\'e group with $U(g)\A(O)U(g)\ad=\A(gO)$.
\item \textbf{Spectrum condition.} The joint spectrum of the translations $U(a)=\ee^{\ii P\cdot a}$ lies in the closed forward light cone.
\item \textbf{Vacuum.} There is a unit vector $\Omega$, unique up to phase, invariant under $U$, and (irreducibility) cyclic for the algebra generated by all $\A(O)$.
\end{enumerate}
Optional structural properties that hold in many models: \emph{Haag duality}\index{Haag's theorem} $\A(O')=\A(O)'$ (for wedges, or double cones in suitable models), \emph{additivity} (a region covered by smaller ones is generated by them), and the \emph{time-slice property} (the algebra of a neighbourhood of a Cauchy surface is the algebra of the whole domain of dependence).

The \emph{quasi-local algebra} $\A=\overline{\bigcup_O\A(O)}^{\norm\cdot}$ is a \cstar-algebra, the vacuum state $\omega_0=\inner\Omega{\,\cdot\,\Omega}$ its Poincar\'e-invariant state, and $(\HH,\pi_0,\Omega)$ the GNS data of $\omega_0$ (\cref{ch:gns}). Other representations (thermal states, charged sectors) arise from other states and are generally inequivalent (\cref{ch:inequiv}).

\begin{keyidea}
A QFT is a net of von Neumann algebras $O\mapsto\A(O)$ with isotony, locality, Poincar\'e covariance, positive energy, and a unique invariant vacuum. All dynamical content lives in how the algebras sit relative to each other.
\end{keyidea}

\subsection{Free-field example}

For the free scalar field of mass $m\ge0$, let $\mathcal S$ be the real symplectic space of Cauchy data and $\mathcal S_O$ the data supported in (the domain of dependence of) $O$. Define
\begin{equation}
\A(O)=\{W(f):f\in\mathcal S_O\}''\subset\Bnd{\mathcal F},
\end{equation}
with $W(f)=\ee^{\ii\phi(f)}$ the Weyl operators on Fock space (\cref{ch:weyl,ch:gns}). Isotony is obvious, locality follows from the Weyl relations because the symplectic form vanishes for spacelike supports, covariance from Poincar\'e invariance of the symplectic form and of the one-particle structure, and the spectrum condition from positivity of the one-particle energy. The Fock vacuum is the vacuum. This is the algebraic version of canonical quantization, defined without reference to a Hamiltonian.

\subsection{Relation to Wightman's framework}

In the Wightman approach one has operator-valued distributions $\phi(f)$ on a dense domain, with the same covariance and spectrum conditions, locality $[\phi(f),\phi(g)]=0$ for spacelike supports, and a cyclic vacuum. From such fields one \emph{builds} nets: $\A(O)$ is the von Neumann algebra generated by bounded functions (e.g.\ spectral projections or exponentials) of the self-adjoint $\phi(f)$ with $\mathrm{supp}f\subset O$, \emph{provided} the commutation of fields at spacelike separation holds in the strong sense (commutation of the spectral projections of the self-adjoint extensions), a technical point that is automatic for free fields and is an assumption in general. Conversely, from a net one cannot generally recover fields (these need additional smoothness), but the net carries all the physical information: scattering theory (Haag--Ruelle), the particle spectrum, and charges (Doplicher--Haag--Roberts, \cref{ch:dhr}).

\begin{whatwrong}
\begin{itemize}
\item The net is a statement about the \emph{vacuum representation}; ``the Hilbert space of the theory'' still means this representation. Haag's theorem (\cref{ch:why_algebras}) says that interacting vacua are not unitarily related to free ones, but \emph{nets} of interacting theories are fine.
\item Gravity does not fit: there is no Poincar\'e group acting on a fixed background and no local algebra of diffeomorphism-invariant observables (\cref{ch:gravity_constraints}). Parts VII--VIII show how the algebra of a region must be redefined.
\item Constructing interacting nets in $d=4$ is an open problem; they exist in $d=2$ and for conformal theories (\cref{ch:typeIII_local}), and free and many integrable models are covered.
\end{itemize}
\end{whatwrong}

\subsection*{Exercises}
\begin{exercise}\label{ex:haagkastler:1}
For the free field show that $[\A(O_1),\A(O_2)]=0$ for spacelike $O_1,O_2$ using the Weyl relation and the vanishing of the commutator function outside the light cone.
\end{exercise}
\begin{exercise}\label{ex:haagkastler:2}
Explain why an algebra $\A(O)$ cannot be type I (i.e.\ $\Bnd{\HH_O}$ for a tensor factor) if $\A(O')=\A(O)'$ and the vacuum is entangled across $\partial O$; anticipate \cref{ch:typeIII_local}.
\end{exercise}
\begin{exercise}\label{ex:haagkastler:3}
In $1+1$ dimensions with a massless free chiral boson, write $\A(I)$ for an interval $I$ of the light ray in terms of $W(f)$, $\mathrm{supp}f\subset I$, and check that translations and dilations act geometrically.
\end{exercise}

\section{Reeh--Schlieder, Bisognano--Wichmann, and the Unruh effect}\label{ch:rs_bw}

\begin{physmotiv}
Two theorems give the vacuum of a QFT its algebraic character. Reeh--Schlieder says that acting with operators from a tiny region on the vacuum produces (approximately) \emph{every} state: the vacuum is cyclic and separating for local algebras, so modular theory applies. Bisognano--Wichmann says what the modular flow is: for a Rindler wedge, it is the Lorentz boost \cite{BisognanoWichmann1975,BisognanoWichmann1976}. Together they make the Unruh effect a theorem about modular flow and show that the abstract Tomita--Takesaki structure of \cref{ch:tomita} is, in QFT, geometry.
\end{physmotiv}

\subsection{Reeh--Schlieder}

\begin{theorem}[Reeh--Schlieder]
Let $O$ be a non-empty open region. Under the axioms of \cref{ch:haagkastler}, the vacuum $\Omega$ is cyclic for $\A(O)$. If moreover the causal complement $O'$ has non-empty interior, $\Omega$ is also separating for $\A(O)$ (it is cyclic for $\A(O')\subset\A(O)'$).
\end{theorem}

\begin{proofsketch}
Let $\Psi\perp\A(O)\Omega$. Consider $F(x_1,\dots,x_n)=\inner\Psi{\phi(x_1)\cdots\phi(x_n)\Omega}$, where translations $x_j-x_{j-1}$ are taken in $O$. By the spectrum condition (positive energy), $F$ extends to an analytic function of the differences $x_j-x_{j-1}$ in the forward tube $\mathrm{Im}\,(x_j-x_{j-1})\in V_+$. It vanishes on an open set (where all points are in $O$), hence everywhere by analyticity (edge-of-the-wedge/unique continuation). So $\Psi$ is orthogonal to the dense set generated by the vacuum, so $\Psi=0$.
\end{proofsketch}

Consequences. (i) \emph{Local operators in an arbitrarily small region can create any state}, including one with a particle at the other side of the universe; no contradiction with causality because the operator is not a local observable of that distant region and it needs a non-local looking combination of unitaries. (ii) No non-zero element of $\A(O)$ annihilates $\Omega$; so $\omega_0$ is a faithful state on $\A(O)$. (iii) Hence we may apply \cref{ch:tomita} to $(\A(O),\Omega)$ for \emph{every} region with non-trivial complement. (iv) Combined with \cref{ch:classification}: $\A(O)$ can have no minimal projection, since that would give a pure normal state incompatible with the analyticity underlying the proof; this is made precise in \cref{ch:typeIII_local}.

\begin{figure}[h]
\centering
\begin{tikzpicture}[scale=0.95]
\fill[physcol!15] (0,0)--(3,3)--(3,-3)--cycle;
\fill[physcol!6] (0,0)--(-3,3)--(-3,-3)--cycle;
\draw[dashed] (-3,-3)--(3,3); \draw[dashed] (-3,3)--(3,-3);
\foreach \r in {0.8,1.5,2.2}{
 \draw[thick,-{Stealth}] plot[domain=-1.0:0.95,smooth,variable=\s] ({\r*cosh(\s)},{\r*sinh(\s)});
 \draw[thick,{Stealth}-] plot[domain=-0.95:1.0,smooth,variable=\s] ({-\r*cosh(\s)},{\r*sinh(\s)});
}
\node[font=\small] at (2.55,-0.2) {$W$};
\node[font=\small] at (-2.55,-0.2) {$W'$};
\draw[->] (-3.4,0)--(3.4,0) node[right,font=\small] {$x$};
\draw[->] (0,-3.4)--(0,3.4) node[above,font=\small] {$t$};
\end{tikzpicture}
\sidecaption{Modular flow for the Rindler wedge $W$ in the Minkowski vacuum (Bisognano--Wichmann). $\sigma_s=\Ad\ee^{-\ii sK}$ with $K=2\pi B$ acts on $\A(W)$ as the Lorentz boost; the orbits are the hyperbolas $x^2-t^2=\mathrm{const}$. On $W'=W^c$ the same $K$ acts as the boost with reversed orientation.}
\end{figure}

\subsection{Bisognano--Wichmann}

Let $W=\{x:x^1>|x^0|\}$ be the right Rindler wedge, $\A(W)$ its algebra, $\A(W')=\A(W)'$ (Haag duality) that of the left wedge. Let $\Lambda_W(s)$ be the boost with rapidity $s$ that preserves $W$ (moving forward in time inside $W$ for $s>0$), $U(\Lambda_W(s))=\ee^{\ii sB}$ with $B$ the boost generator.

\begin{theorem}[Bisognano--Wichmann]\label{thm:BW}
For the vacuum and a wedge algebra in a Wightman theory,
\begin{equation}
\Delta^{\ii t}_{\Omega}=U\big(\Lambda_W(-2\pi t)\big),\qquad J_\Omega=\Theta\,U(R_W(\pi)),
\end{equation}
where $\Theta$ is the CPT operator and $R_W(\pi)$ the rotation by $\pi$ in the plane orthogonal to the $(x^0,x^1)$ plane (so $J$ maps $\A(W)$ to $\A(W')$). In particular the modular Hamiltonian is $K=-\log\Delta=2\pi B$.
\end{theorem}
(This is a theorem for theories obeying the Wightman axioms; the result was proven in the 1970s by Bisognano and Wichmann, and extended to nets by various authors, so that $\Delta^{it}$ acts geometrically on the whole net of the wedge and its subregions.)

\begin{proofsketch}
Formally: the vacuum correlation functions are analytic in the forward tube. Rotating the Euclidean angle by $\pi$ in the $(x^0,x^1)$-plane maps the right wedge to the left wedge: $\Theta U(R(\pi))$ is the antilinear map $a\Omega\mapsto a\ad\Omega$, i.e.\ $S=J\Delta^{1/2}$ is the analytic continuation of the boost to ``imaginary rapidity $\ii\pi$'', so $\Delta^{1/2}=\ee^{-\pi B}$, which is $\ee^{-K/2}$ with $K=2\pi B$. Rigorous proofs use the spectrum condition to continue $U(\Lambda(s))$ analytically into the strip $0<\mathrm{Im}\,s<\pi$ on vectors $a\Omega$.
\end{proofsketch}

\begin{keyidea}
The modular flow of the Rindler wedge in the vacuum is the boost, with $K=2\pi B$ and $J$ the CRT reflection. The abstract Tomita--Takesaki objects $(\Delta,J)$ are Lorentz geometry.
\end{keyidea}

\subsection{The Unruh effect as modular flow}

Combining \cref{thm:BW} with the KMS property of \cref{ch:tomita}: the vacuum restricted to $\A(W)$ is KMS for the boost group at inverse temperature $\beta=2\pi$ (\cref{ch:kms}). An observer with uniform acceleration $a$ has trajectory $x(\tau)=(a^{-1}\sinh a\tau,a^{-1}\cosh a\tau)$ in $W$, and proper time $\tau$ corresponds to rapidity $s=a\tau$. The proper-time evolution is then $\alpha_\tau=\Ad U(\Lambda_W(a\tau))$, which is KMS at $\beta_\tau=2\pi/a$. A detector coupled to the field along this worldline therefore sees the vacuum as a thermal state at
\begin{equation}
T_U=\frac{a}{2\pi}\quad(\hbar=c=k_B=1).
\end{equation}
Here the temperature is a property of the vacuum's modular structure and not of any particle content: the vacuum \emph{is} an equilibrium state of the boost dynamics, in a representation (the vacuum Fock space) in which $\A(W)$ has no density matrix. The same logic gives the Hawking temperature as the KMS temperature at a Killing horizon (\cref{ch:curved_hawking}) and the Gibbons--Hawking temperature in the de Sitter static patch (\cref{ch:static_patch}).

\begin{whatwrong}
``The accelerated observer sees particles'' is a statement about the choice of the \emph{state restricted to the wedge}; the observer's algebra is $\A(W)$ in both descriptions. The inequivalence of the thermal Rindler and the Minkowski particle notions is the Fock-space inequivalence of \cref{ch:weyl}: Rindler quanta are not a unitary rotation of Minkowski quanta. Equally, ``the Unruh effect needs a horizon'' is built in: $\A(W)$ is a type III algebra \emph{because} the wedge has an entangling surface.
\end{whatwrong}

\subsection*{Exercises}
\begin{exercise}\label{ex:rs_bw:1}
For the massless scalar in $1+1$ dimensions (work with the chiral currents $\partial_\pm\phi$ to avoid infrared problems), the right wedge algebra is generated by right- and left-moving chiral fields on half-lines. Show that on the one-particle space the boost acts as a dilation $p\mapsto\ee^{s}p$ of the null momentum and that $\Delta^{1/2}=\ee^{-\pi B}$ maps $p>0$ momenta as an analytic continuation to $\ee^{\ii\pi}p$.
\end{exercise}
\begin{exercise}\label{ex:rs_bw:2}
Using $K=2\pi B$ show that $\langle\Omega|K|\Omega\rangle=0$, while for $\Psi=u\Omega$ with $u\in\A(W)$ unitary and $\Psi\ne\Omega$ one has $\langle\Psi|K|\Psi\rangle=S_{\rm rel}(\Psi\|\Omega)>0$ (\cref{ch:modham}). (Positivity of the boost energy density pointwise is \emph{not} available for general states, so this is the correct route.)
\end{exercise}
\begin{exercise}\label{ex:rs_bw:3}
Derive $T_U=a/2\pi$ from $\beta=2\pi$ in rapidity and $s=a\tau$. Restore $\hbar,c,k_B$.
\end{exercise}
\begin{exercise}\label{ex:rs_bw:4}
Argue from Reeh--Schlieder that the vacuum state restricted to $\A(O)$ is faithful and conclude it is not pure: a pure normal state of a von Neumann algebra forces a minimal projection, so $\A(O)$ would be type I, contradicting the type III property of \cref{ch:typeIII_local}.
\end{exercise}

\section{Local algebras are type III$_1$}\label{ch:typeIII_local}

\begin{physmotiv}
We now collect the evidence for the claim on which the whole book rests: algebras of local regions in QFT are of type III$_1$. This explains in one stroke the divergence of entanglement entropy, the absence of local pure states, and the Reeh--Schlieder property, and tells us exactly what must change when gravity is included: we need a trace.
\end{physmotiv}

\subsection{The statement}

\begin{center}
\emph{In a (reasonable) local QFT, the algebra $\A(O)$ of a bounded region with non-trivial causal complement is the hyperfinite factor of type III$_1$.}
\end{center}

Parts of this are theorems; other parts require assumptions. We list the logic, separating what is proven under general axioms from what needs more.

\begin{enumerate}
\item \textbf{Cyclic and separating vacuum} (Reeh--Schlieder, \cref{ch:rs_bw}) $\Rightarrow$ modular theory applies, and $\omega_0|_{\A(O)}$ is a faithful normal state.
\item \textbf{No minimal projections, no trace.} A type I factor has pure normal states. Pure normal states are incompatible with analyticity in the translations (the energy-momentum spectrum condition), so $\A(O)$ is not type I. Under conditions that include the \emph{split property}\index{split property} and scaling-type assumptions one shows that no finite trace exists; local algebras are of type III (Fredenhagen, Buchholz, and others; see the literature).
\item \textbf{Modular spectrum.} For wedges, $\Delta^{it}$ is the boost and $\spec\Delta=[0,\infty)$ because the boost generator has spectrum $\R$ (\cref{ch:connes_class}, exercise). Together with factoriality this shows the wedge algebra is type III$_1$ (Driessler; Longo; Fredenhagen; see Buchholz--D'Antoni--Longo for the passage to double cones).
\item \textbf{Uniqueness.} If the net satisfies nuclearity/split conditions and has a non-trivial scaling limit, $\A(O)$ is the unique \emph{hyperfinite}\index{hyperfinite factor} type III$_1$ factor $R_\infty$ (Haagerup's theorem \cite{Haagerup1987} identifies all injective type III$_1$ factors; Buchholz--D'Antoni--Fredenhagen \cite{BDF1987} obtain the universal type III$_1$ structure from a scaling limit, and Buchholz--D'Antoni--Longo \cite{BuchholzDAntoniLongo1990} treat the nuclearity/split side; the hypotheses differ between these papers, so consult them before quoting a precise statement).
\end{enumerate}

The upshot is that, for the theories we care about (free fields, conformal theories, and any theory with a good scaling limit in the UV), local algebras are \emph{all isomorphic} to one single abstract algebra: $R_\infty$, the Araki--Woods factor of \cref{ch:connes_class}. Differences between theories are not in the isomorphism class of $\A(O)$ but in how the family of algebras sits in $\Bnd{\HH}$.

\subsection{Consequences}

\begin{center}
\renewcommand{\arraystretch}{1.3}
\resizebox{\ifdim\width>\linewidth\linewidth\else\width\fi}{!}{%
\begin{tabular}{@{}p{4.6cm}p{8cm}@{}}
\toprule
Property of type III$_1$ & Physics\\
\midrule
No trace; no density matrix & Region entropy $-\Tr\rho\log\rho$ ill defined; UV divergence $\sim\mathrm{Area}/\epsilon^{d-2}$\\
No pure normal states & No local state is pure; any state is entangled with the environment\\
All projections $e\sim\id$ & Local operators can generate any state from the vacuum (Reeh--Schlieder)\\
Modular spectrum $[0,\infty)$; flow outer & The boost acts as an outer automorphism: Unruh/Hawking temperature without a Hamiltonian in the region\\
Maximal Bell violation & Entanglement is as strong as quantum mechanics allows (\cref{ch:subsystems})\\
Entanglement dense & Product states do not exist as normal states of $\A(R)\vee\A(L)$\\
\bottomrule
\end{tabular}}
\end{center}

\begin{whatwrong}
``Type III$_1$'' is an exact statement about the continuum algebra. The lattice, with finite-dimensional local Hilbert spaces, is type I, and the area-law entropy there is finite and cutoff dependent. The continuum limit loses the factorization, not just the finiteness. Equally, the absence of a trace does \emph{not} mean entropy differences are meaningless: relative entropy and the first law survive (\cref{ch:modham}).
\end{whatwrong}

\subsection{What gravity must supply}

The structure theory of \cref{ch:connes_class,ch:crossed,ch:takesaki} tells us the minimal modification that yields a well-defined entropy: take the crossed product by the modular flow, obtaining a type II$_\infty$ algebra with a trace. The remaining question, of \emph{physics}, is why gravity should do this. The answer (Part VII) is that the modular flow of a horizon is a symmetry whose generator is a \emph{constraint} in a diffeomorphism-invariant theory: gauge-invariant observables must commute with it, which means they must be dressed to a physical clock, and the algebra of dressed observables is the crossed product.

\begin{keyidea}
Local algebras of QFT are (in nice theories) the hyperfinite III$_1$ factor. Everything about vacuum entanglement that is divergent or non-factorizing is a symptom of this. The minimal cure is the crossed product, and gravity is expected to provide it.
\end{keyidea}

\subsection*{Exercises}
\begin{exercise}\label{ex:typeIII_local:1}
Show that a type I$_\infty$ factor $\Bnd K\otimes\id\subset\Bnd{K\otimes K'}$ would give $\A(O)$ a pure normal state, and explain why this contradicts the fact that the vacuum is cyclic for $\A(O')=\A(O)'$.
\end{exercise}
\begin{exercise}\label{ex:typeIII_local:2}
Starting from the free-field lattice entropy $S=\frac13\log(L/\epsilon)$ in $1+1$ dimensions, explain why the $\epsilon$-dependence cannot be removed by any state-independent subtraction that is covariant in the continuum, in the language of type III.
\end{exercise}
\begin{exercise}\label{ex:typeIII_local:3}
List which of the following are type III$_1$: the wedge algebra of a free field; the algebra of a double cone in a CFT; the algebra of the whole Minkowski space in the vacuum representation; the algebra of an ordinary quantum mechanical particle. Justify each.
\end{exercise}

\section{Superselection and DHR}\label{ch:dhr}
\skippable

\begin{physmotiv}
How do electric charge, baryon number and spin-statistics show up in a theory formulated only through local observables? Doplicher, Haag and Roberts \cite{DHR1971,DHR1974} showed that charged sectors are \emph{other representations} of the same quasi-local algebra (\cref{ch:inequiv}), distinguished by being indistinguishable from the vacuum outside a bounded region. The set of such sectors forms a category with products (fusion) and a braiding. Statistics (Bose/Fermi, or parastatistics, or anyons in $2+1$ dimensions) is read off from the braiding, and in $d\ge3+1$ the gauge group is recovered (Doplicher--Roberts). This section is a conceptual overview; it is not needed for the gravity chapters.
\end{physmotiv}

\subsection{DHR sectors}

Let $\pi_0$ be the vacuum representation of the quasi-local algebra $\A$. A representation $\pi$ satisfies the \emph{DHR selection criterion} if for every double cone $O$,
\begin{equation}
\pi\restriction\A(O')\ \cong\ \pi_0\restriction\A(O')\quad\text{(unitarily equivalent)}.
\end{equation}
This says the charge is localized in a bounded region: far from it, $\pi$ looks like the vacuum. Such $\pi$ can be realized as $\pi_0\circ\rho$ with $\rho$ a \emph{localized endomorphism} of $\A$ ($\rho=\id$ on $\A(O')$). Equivalence classes of irreducible such $\rho$ are the \emph{superselection sectors}\index{superselection sector}.

\subsection{Structure}
\begin{itemize}
\item \textbf{Composition.} $\rho_1\circ\rho_2$ is again localized: charges can be fused. The sectors form a (tensor) category $\mathrm{DHR}(\A)$ with direct sums, conjugates (antiparticles, for finite statistics) and the vacuum as unit.
\item \textbf{Statistics.} Localizing $\rho_1,\rho_2$ at spacelike separation, the two compositions $\rho_1\rho_2$ and $\rho_2\rho_1$ are unitarily equivalent by a \emph{statistics operator} $\varepsilon(\rho_1,\rho_2)$. In $d\ge3+1$, $\varepsilon^2=1$ on irreducible sectors: Bose or Fermi (or parastatistics of finite order, then resolved by the next theorem). In $2+1$ dimensions $\varepsilon$ generates a braid group representation: anyons.
\item \textbf{Dimension.} Every sector has a dimension $d(\rho)\ge1$, $d(\rho)=1$ iff abelian charge; $d(\rho)$ is the \emph{quantum dimension} and determines how a sector contributes to entropy (for instance, topological entanglement entropy).
\end{itemize}

\begin{theorem}[Doplicher--Roberts, $d\ge3+1$, conceptual statement]
Under the above conditions there is a compact group $G$ and a net of ``field'' algebras $\mathcal F(O)$ with graded locality, carrying a $G$-action whose fixed points are $\A(O)$, and whose sectors are exactly the representations of $G$. The gauge group is reconstructed from the observables.
\end{theorem}

In this sense gauge symmetry is not an input, it is a derived property of a local net. It also fits with the central-decomposition viewpoint of \cref{ch:factors}: superselection sectors are labels of a decomposition (a generalized centre) of the algebra of \emph{all} (including charged) operators.

\begin{whatwrong}
DHR theory assumes a mass gap or at least Haag duality and a localization of the charge in a \emph{bounded} region; theories with long-range forces (electric charge with massless photons, Gauss's law at infinity, infraparticles) require a modified analysis (Buchholz--Fredenhagen). And, in gravity, \cref{ch:gravity_constraints} will argue there are no local charges and no non-trivial sectors localized in bounded regions.
\end{whatwrong}

\subsection*{Exercises}
\begin{exercise}\label{ex:dhr:1}
Show that a localized endomorphism of the Weyl algebra of a free field, $\rho(W(f))=\ee^{\ii\lambda\,\ell(f)}W(f)$ with $\ell$ a symplectic functional supported in $O$, is DHR-localized, and describe its charge (a coherent shift).
\end{exercise}
\begin{exercise}\label{ex:dhr:2}
Using the graded commutation of fermionic operators, show that the fermion parity $(-1)^F$ is a $\Z_2$ gauge group whose representations give two sectors (even/odd).
\end{exercise}

